\section*{الفصل \ref{ch:b3:surfaces-catalysis} --- السطوح والامتزاز والحفز غير المتجانس}

\begin{solution}{exo:b3:surfaces-catalysis:1}
$K = \theta/((1 - \theta)p) = 0.40/(0.60 \times 2.0) = \qty{0.33}{kPa^{-1}}$؛ وعند \qty{10}{kPa}، $\theta = 3.33/4.33 = 0.77$.
\end{solution}

\begin{solution}{exo:b3:surfaces-catalysis:2}
$-20$: \omterm{def:b3:surfaces-catalysis:physi-chemi}{امتزاز فيزيائي} (من رتبة أنتالبي التكاثف)؛ $-150$: \omterm{def:b3:surfaces-catalysis:physi-chemi}{امتزاز كيميائي}، وهو
وحده يمكن أن يكون تفككيًا.
\end{solution}

\begin{solution}{exo:b3:surfaces-catalysis:3}
(100): موقع علوي واحد، وموقعان جسريان (أربع روابط تتقاسم كلًّا منها ذرتان)، وتجويف رباعي
واحد (أربعة تجاويف يتقاسم كلًّا منها أربع ذرات). (111): موقع علوي واحد، وثلاثة مواقع جسرية، وتجويفان
ثلاثيان (ستة مثلثات حول كل ذرة، يتقاسم كلًّا منها ثلاث ذرات).
\end{solution}

\begin{solution}{exo:b3:surfaces-catalysis:4}
$\num{3.0e-6}/\num{1.5e-5} = \qty{0.20}{s^{-1}}$.
\end{solution}

\begin{solution}{exo:b3:surfaces-catalysis:5}
$\theta/(1 - \theta) = (Kp)^{1/2}$: القيمة $0.111$ عند \qty{1.0}{kPa}، تتضاعف عند \qty{4.0}{kPa}: $0.222$، ومنه $\theta = 0.18$.
\end{solution}

\begin{solution}{exo:b3:surfaces-catalysis:6}
المقام $1 + 10 + 2.5 = 13.5$: $\theta_{\ce{CO}} = 0.74$، $\theta_{\ce{O2}} = 0.19$، والمواقع الحرة 0.07. 
وسرعة لانغموير--هينشلوود، المتناسبة مع $\theta_{\ce{CO}}\theta_{\ce{O2}}$، يحدّها
شحّ الأكسجين على سطح مزدحم بالجزيء CO.
\end{solution}

\begin{solution}{exo:b3:surfaces-catalysis:7}
$\Delta_{\mathrm{ads}}H = -R\ln(8.0/1.0)/(1/300 - 1/350) = -8.314 \times 2.079/\num{4.76e-4} = \qty{-36}{kJ/mol}$.
\end{solution}

\begin{solution}{exo:b3:surfaces-catalysis:8}
$n_{\mathrm m} = 120/22\,414 = \qty{5.35e-3}{mol}$؛ والمساحة $= \num{5.35e-3} \times \num{6.022e23} \times \num{0.162e-18} =
\qty{522}{m^2.g^{-1}}$.
\end{solution}

\begin{solution}{exo:b3:surfaces-catalysis:9}
لانغموير--هينشلوود: يجب أن يكون المتفاعلان كلاهما ممتزين على مواقع متجاورة؛
ويرتبط CO بقوة، فلا يترك عند الضغط المرتفع مكانًا للأكسجين.
\end{solution}

\begin{solution}{exo:b3:surfaces-catalysis:10}
عند تغطية ضعيفة: $100 - 60 = \qty{40}{kJ/mol}$؛ وعند تغطية مرتفعة: \qty{100}{kJ/mol}. فمنحنى $\ln r$
بدلالة $1/T$ حاد الانحدار (الميل $-100/R$) عند درجة الحرارة المنخفضة، حيث يكون السطح
ممتلئًا، وأقل انحدارًا (الميل $-40/R$) عند درجة الحرارة المرتفعة، حيث يفرغ: منحنى
ينثني بين الحالتين.
\end{solution}

\begin{solution}{exo:b3:surfaces-catalysis:11}
$4 \times 0.10 = 0.40$ من المواقع مسدودة: فيبقى 60\,\% من النشاط من أجل موقع
واحد، ونحو $0.60^2 = 36\,\%$ من أجل موقعين حرين متجاورين.
\end{solution}

\begin{solution}{exo:b3:surfaces-catalysis:12}
الأوسط (\qty{-250}{kJ/mol}): فالأول يربط النيتروجين بضعف مفرط فلا يكسر \ce{N2}
بسهولة، والثالث بقوة مفرطة، فتبقى ذرات النيتروجين على السطح
وتسده.
\end{solution}

\begin{solution}{pb:b3:surfaces-catalysis:1}
\textbf{1.} يتكاثف النيتروجين قرب \qty{77}{K} تحت الضغط الجوي: فيُبلغ مجال كامل من
الضغوط النسبية، وتتكوّن طبقات متعددة.
\textbf{2.} دون 0.05 يخالف عدمُ تجانس السطح فرضياتِ BET، وفوق 0.30 يخالفها التكاثفُ في
المسامّ.
\textbf{3.} \num{1.278e-3}، \num{2.410e-3}، \num{3.453e-3}، \num{4.621e-3}، \num{5.659e-3}، \num{6.813e-3} \unit{g.cm^{-3}}.
\textbf{4.} $v_{\mathrm m} = 1/(0.02205 + 0.000180) = \qty{45.0}{cm^3.g^{-1}}$.
\textbf{5.} $c = 1 + 0.02205/0.000180 = 124$: فالطبقة الأولى ترتبط أقوى بكثير مما
يفعل السائل.
\textbf{6.} إنه ضئيل جدًا مقارنة بجداء الميل في $x$، فخطؤه النسبي كبير؛
لكن $v_{\mathrm m}$ تتوقف على مجموع الميل والتقاطع، الذي يغلب عليه الميل.
\textbf{7.} $45.0/22\,414 = \qty{2.01e-3}{mol.g^{-1}}$.
\textbf{8.} $\num{2.01e-3} \times \num{6.022e23} \times \num{0.162e-18} = \qty{196}{m^2.g^{-1}}$.
\textbf{9.} حامل الألومينا: فالبلاتين المكشوف (الجزء الثالث) مساحته نحو \qty{0.9}{m^2.g^{-1}}.
\textbf{10.} المسامّ التي عرضها بحجم الجزيئات تمتلئ كليًا عند ضغط منخفض جدًا، لا طبقةً
بعد طبقة، فلا يمكن تحديد أي \omterm{def:b3:surfaces-catalysis:coverage}{طبقة أحادية}.
\textbf{11.} $2 \times 0.205/22\,414 = \qty{1.83e-5}{mol.g^{-1}}$.
\textbf{12.} $0.0100/195.08 = \qty{5.13e-5}{mol.g^{-1}}$.
\textbf{13.} $D = 1.83/5.13 = 0.357$.
\textbf{14.} $a^3/4 = (0.3923)^3/4 = \qty{0.01509}{nm^3}$.
\textbf{15.} $\frac{\sqrt3}{4}(0.3923)^2 = \qty{0.0666}{nm^2}$، وهو الوجه الأكثف؛ ويعطي (100) و(110) القيمتين 0.0770 
و\qty{0.1088}{nm^2}، ويعطي متوسط الكثافات السطحية الثلاث للمواقع \qty{0.0807}{nm^2}.
\textbf{16.} تحوي الكرة $\pi d^3/6v_{\mathrm{at}}$ ذرة، منها $\pi d^2/a_{\mathrm{at}}$ على السطح:
$D = 6v_{\mathrm{at}}/(a_{\mathrm{at}}d)$.
\textbf{17.} $d = 6 \times 0.01509/(0.0807 \times 0.357) = \qty{3.1}{nm}$.
\textbf{18.} ذرة H واحدة لكل ذرة Pt سطحية؛ وجسيمات كروية من حجم واحد (فالنتيجة
متوسط مرجَّح بالسطح)؛ وكل البلاتين مختزل ومتاح؛ ولا انسكاب للهيدروجين
على الحامل.
\textbf{19.} $\num{2.0e-5}/\num{1.83e-5} = \qty{1.1}{s^{-1}}$.
\textbf{20.} $\num{2.0e-5}/0.0100 = \qty{2.0e-3}{mol.g^{-1}.s^{-1}}$ لكل غرام من البلاتين.
\textbf{21.} $D = 1.12/10 = 0.112$ بدل 0.357: أي بعامل 3.2.
\textbf{22.} أن التفاعل حساس للبنية: فمواقعه النشطة (الحواف،
والرؤوس، وأوجه معيّنة) ليست نسبة ثابتة من السطح.
\textbf{23.} يعدّ الأحداث لكل \omterm{def:b3:complex-kinetics:enzyme}{موقع نشط}، فيزيل أثر مقدار
الفلز المكشوف.
\textbf{24.} \textbf{متوسط قطر جسيمات البلاتين $\approx \qty{3.1}{nm}$.}
\end{solution}
